% ════════════════════════════════════════════════════════════════════
%  front_pages/abstract_nepali.tex
%  Nepali Abstract (शोधसार) — GNN EGFR Thesis
%  Compiled with XeLaTeX
% ════════════════════════════════════════════════════════════════════

\cleardoublepage

\begin{center}
  {\fontsize{18}{22}\selectfont\bfseries{\nepalifont शोधसार}}
\end{center}
\addcontentsline{toc}{chapter}{\texorpdfstring{{\nepalifont शोधसार}}{Shodhasar}}
\thispagestyle{plain}

\vspace{1cm}

{\fontsize{12}{18}\selectfont

% ── Sense-summary (≤150 words per IoST Regulation 2082, Appendix H) ───
{\nepalifont फोक्सोको क्यान्सरमध्ये सबैभन्दा बढी पाइने नन-स्मल सेल
लङ क्यान्सर}
(NSCLC)
{\nepalifont मा}
EGFR
{\nepalifont उत्परिवर्तन प्रमुख कारक हो, तर अवरोधकहरूप्रतिको प्रतिरोध
र उपचारको उच्च लागत ठूलो चुनौती बनेको छ। यस अध्ययनले}
EGFR
{\nepalifont अवरोधकको वर्गीकरण र नेपाली औषधीय वनस्पतिका यौगिकहरूको
भर्चुअल स्क्रिनिङका लागि ग्राफ न्युरल नेटवर्क}
(GNN)
{\nepalifont ढाँचा प्रस्तुत गर्दछ।}
ChEMBL
{\nepalifont बाट लिइएका १०,५२३ यौगिकमा}
Bemis-Murcko
{\nepalifont स्क्याफोल्ड विभाजन प्रयोग गरी पाँच}
GNN
{\nepalifont र दुई आधारभूत मोडेल गरी सात मोडेलको तुलना गरियो।
फिंगरप्रिन्ट-आधारित}
Random Forest
(AUROC 0.8827)
{\nepalifont ले सर्वोत्तम प्रदर्शन गर्दै सबै}
GNN
{\nepalifont लाई उछिन्यो। दुई उत्कृष्ट मोडेलद्वारा सात नेपाली
वनस्पतिका ५२७ नयाँ यौगिकको स्क्रिनिङ गर्दा कुनै पनि यौगिकले}
p $\geq$ 0.80
{\nepalifont पार गरेन;}
xanthone isogentisin
{\nepalifont र}
flavone luteolin
{\nepalifont मध्यम-विश्वासका सम्भावित यौगिक रहे। यसले}
GNN
{\nepalifont-आधारित विधिले नेपाली परम्परागत वनस्पतिबाट औषधि-योग्य
सम्भावित यौगिक पहिचानमा सहयोग गर्न सक्ने देखाउँछ।}

} % end fontsize

\vspace{0.8cm}

% ── Keywords (five, Font 10 italic per Appendix H) ───────────────────
\noindent
{\fontsize{10}{12}\selectfont
{\nepalifont\textbf{मुख्य शब्दहरू:}}
\textit{EGFR {\nepalifont अवरोधक}, {\nepalifont ग्राफ न्युरल नेटवर्क},
{\nepalifont भर्चुअल स्क्रिनिङ}, {\nepalifont नेपाली औषधीय वनस्पति},
{\nepalifont औषधि खोज}}}